% Pursuing Stacks (cotes 134-2 to 134-8): archive cote 134-5, site page 126 (PDF page 127), typescript page 313.
% typescript. Draft compared: ps5.tex:1051-1108.
% AI-transcribed by Claude Opus 5.5 (claude-opus-5-5, adaptive thinking, effort high), Message Batches API, batch ps3 (msgbatch_01XdFmuVVBAcgc7EFxdtCqsj), 2026-11-12; not yet checked by eye.
\scanpage{134-5/126}
\mspage{313}
to \unsure{ismmorphism|isomorphism})  as the homotopy type of a \uline{space admitting a finite}
\uline{triangulation}.  In terms of \unsure{ss|Ss} sets i.e. \inserted{of} the modelizer \hand{$\Delta$}$^\wedge$,  the natu-%
ral finiteness condition,  suggested by the algebraic formalism,  is
that the homotopy type be isomorphic to one defined by an object ``\uline{of}
\uline{finite presentation}'' \uline{in} \hand{$\Delta$}$^\wedge$ ,  namely one which is a direct limit of
a \uline{finite} diagram in \hand{$\Delta$}$^\wedge$,  made up with simplices  i.e. coming from a
diagram in \hand{$\Delta$}.  It is clear that the first finiteness condition implies
the second,  by using a total order on the set of vertices of the trian-%
gulation.  T\hand{h}e converse  should'nt be hard either,  using an induction
argument on the number of simplices occuring in the diagram,  and using
the fact that any \xout{subobject} quotient object in \hand{$\Delta$}$^\wedge$ of a simplex is
again a simplex,  hence also a subobject of a simplex is a union of
subsimplices;  from this should follow by induction that the geometric
realization of a  ss set of finite presentation is endowed with a
natural \inserted{compact} \uline{piecewise} \uline{linear structure},  and hence can be \inserted{finitely} triangulated. Pre-%
sumably,  one can even find a canonical triangulation,  using twofold
barycentric subdivision Sd$^2$ (again !)  on any simplex.  All this is su-%
rely standard knowledge,  and I do'nt feel like diving into technicali-%
ties on this matter,  unless I am forced to.

If we start with an arbitrary test category A,  the notion of an
object of finite presentation in A$^\wedge$ still makes sense.  Indeed,  in any
category M,  stable \deleted{under} \inserted[typed]{filtering} small direct limits,  we may define object\hand{s} of
finite presentation as those for which the corresponding \uline{covariant}
functor
\[ Y \hand{$\longmapsto$} \mathrm{Hom}(X,Y) \]
commutes to filt\xout{r}ering  direct limits.  If M is stable under \xout{anyxtyp}
some type of finite direct limits,  say \inserted{under} any finite direct limits,  then
so is the \inserted[typed]{full} subcategory $M_{fp}$  of objects of finite presentation of M .
In the case when M=A$^\wedge$,  A any small category,  it is obvious that ob-%
jects of A,  and hence finite direct limits of such,  are of finite pre-%
sentation,  and it is not difficult to show that the converse equally
holds.  As a matter of fact,  A$^\wedge_{fp}$  \uline{can be viewed as the solution of the}
2-\uline{universal} \uline{problem of} ``\uline{adding finite direct limits to A}'' .

Coming back to the case when A is a test category\hand{,} and hence A$^\wedge$ is
modelizing,  one may ask for conditions \inserted{upon A} which ensure that the homotopy
types of finite presentation are exactly those isomorphic to the homoto-%
py types defined by objects of A$^\wedge_{fp}$ .  One \deleted{is clear} \inserted{expects} that some stringent

\margin{126}
